% !TEX root = ../metatime_monograph.tex
\chapter{Cross-Relation with the Secret-of-a-Half Formalism}
\label{app:secret-half-cross-relation}

\claimstatus{Cross-framework formal interface.  Exact statements are separated
from TIR postulates and open physical or zeta-function bridges.}

% CITATION_CONTEXT_V10_9_BEGIN
\paragraph{Established context.}
The binary-entropy statement used below is standard Shannon information theory
\cite{shannon1948}, while the local statistical geometry of Bernoulli families
belongs to the Fisher--Rao framework \cite{Fisher1925,Rao1945}.  Reciprocal
inversion and complement conjugacy in projective odds are elementary algebraic
facts.  The particular cross-identification with the Metatime normalization and
the DHSE/TIR interface are constructions of the present research programme.
% CITATION_CONTEXT_V10_9_END

\section{Purpose and non-circularity rule}

The sibling \emph{Secret of a Half} programme and Metatime/TIR share the
quantities $1/2$ and $\ln2$, but they use them in different logical roles.  A
cross-relation is useful only if the direction of implication is explicit and
neither framework is used to prove one of its own assumptions through the
other.

The interface adopted here is therefore typed:
\[
\boxed{
\text{exact half-side theorem}
\;\longrightarrow\;
\text{TIR structural definition}
\;\longrightarrow\;
\text{exact conditional consequence}
}
\]
with open physical interpretations kept outside that exact chain.

\section{Exact half-side input}

For the Bernoulli binary entropy
\[
H_2(p)=-p\ln p-(1-p)\ln(1-p),
\qquad p\in(0,1),
\]
one has the standard exact result
\[
H_2'(p)=\ln\frac{1-p}{p},
\qquad
H_2''(p)=-\frac1p-\frac1{1-p}<0.
\]
Hence
\begin{equation}
\boxed{
\operatorname*{arg\,max}_{0<p<1} H_2(p)=\frac12,
\qquad
H_2\!\left(\frac12\right)=\ln2
}.
\label{eq:half-entropy-interface}
\end{equation}

The same binary midpoint is also the unique fixed point of complement,
\[
p\mapsto1-p,
\]
and, under the projective odds coordinate
\[
q=\frac{p}{1-p},
\]
it corresponds to the unique positive fixed point $q=1$ of reciprocal inversion
$q\mapsto1/q$.  These fixed-point statements are exact and do not require the
Metatime normalization.

\section{Insertion into the Metatime normalization}

TIR separately defines
\begin{equation}
\kappa\equiv\frac{\ln2}{24\pi}.
\label{eq:half-interface-kappa}
\end{equation}
Equation~\eqref{eq:half-entropy-interface} identifies the numerator with the
maximum binary entropy, but it does \emph{not} derive the denominator $24\pi$.
The discrete orientation/tetrahedral/reference-phase normalization remains the
TIR model postulate documented in Appendix~\ref{app:kappa}.

Consequently the valid logical chain is
\begin{equation}
\boxed{
\frac12
\xrightarrow{\;H_2\;}
\ln2
\xrightarrow{\;\text{TIR definition}\;}
\kappa
}
\label{eq:half-to-kappa-chain}
\end{equation}
where the first arrow is exact information theory and the second arrow is a
model-definition step.

\section{Phase-rate closure}

For angular phase rate
\[
\omega=\frac{d\phi}{dt}=2\pi f
\]
and the TIR definition
\[
d\mathcal I=\kappa\,d\phi,
\]
Appendix~\ref{app:kappa} gives
\begin{equation}
\boxed{
\Gamma_{\mathcal I}
=\frac{d\mathcal I}{dt}
=\kappa\omega
=\frac{\ln2}{12}f
}.
\label{eq:half-interface-phase-rate}
\end{equation}
Thus the cross-framework chain may be written compactly as
\begin{equation}
\boxed{
\frac12
\longrightarrow
\ln2
\longrightarrow
\frac{\ln2}{24\pi}
\longrightarrow
\frac{\ln2}{12}f
}
\end{equation}
provided the logical types of the arrows are retained.

Equivalently, a complete angular cycle carries
\[
\Delta\mathcal I_{\rm cycle}=2\pi\kappa=\frac{\ln2}{12}.
\]
The cancellation of $\pi$ is therefore a consequence of converting an angular
rate to cyclic frequency; it is not an independent evaluation of $\pi$.

\section{Constraint count}

For
\[
\mathbf q=(\kappa,\omega,f,\Gamma_{\mathcal I})
\]
the three constraints
\[
\kappa-\frac{\ln2}{24\pi}=0,
\qquad
\omega-2\pi f=0,
\qquad
\Gamma_{\mathcal I}-\kappa\omega=0
\]
have rank three.  Conditional on the TIR definitions, the subsystem therefore
has one continuous degree of freedom and is globally parametrized by $f$:
\[
\mathbf q(f)=
\left(
\frac{\ln2}{24\pi},
2\pi f,
f,
\frac{\ln2}{12}f
\right).
\]
This is a parameter-counting theorem inside the declared subsystem; it does not
supply an empirical value of $f$.

\section{A necessary negative result: self-duality is not extremality}

The half formalism also supplies an important boundary condition for future TIR
uses of self-duality.  In the finite DHSE-001 Stage-M M\"obius universe, the
forcing-count function obeys reciprocal symmetry
\[
N_n(q)=N_n(1/q),
\]
yet the self-dual point $q=1$ is the unique global maximum only for declared word
lengths $n=2,3$.  At $n=1$ and $n=4$, the exact maxima split into reciprocal
off-centre components.

Therefore
\begin{equation}
\boxed{
\text{self-duality or reciprocal symmetry}
\;\not\Rightarrow\;
\text{dynamical extremum at the fixed point}
}
\label{eq:self-duality-not-extremality}
\end{equation}
without an additional condition such as positivity, convexity, monotonicity,
variational closure, or another explicit theorem.

This negative theorem is relevant to TIR whenever a future construction attempts
to infer physical preference, stability, or attractor status from a self-dual
coordinate alone.

\section{Claim boundary}

The cross-relation established in this appendix does not prove any of the
following:
\begin{enumerate}
\item that standard quantum mechanics independently derives
      $\kappa=\ln2/(24\pi)$;
\item that $\Gamma_{\mathcal I}$ is a measured physical surface-refresh rate;
\item that reciprocal self-duality alone selects a dynamical attractor;
\item that the current Secret-of-a-Half construction closes a canonical bridge
      from non-trivial zeta zeros to the critical line;
\item that the phase-rate identity by itself validates any particle-sector mass
      assignment.
\end{enumerate}

The established interface is narrower and stronger: the exact binary midpoint
produces $\ln2$; TIR uses that quantity in a declared structural normalization;
and the resulting phase-rate and constraint-manifold statements then follow
exactly conditional on those definitions.
